\emph{Bilocal composite: exact reduction of the two-point-kernel gap and conditional non-degeneracy bounds.} This paragraph records derivations that discharge the two-point-kernel non-degeneracy step under explicit hypotheses and scope what the bilocal reading of Assumption~\ref{ass:nondeg} can supply. Let \(M(x,y)=\Phi(x)\,W(x,y)\,\Phi(y)^\dagger\) be the bilocal composite with transport coefficients \(W(x,y)\in\mathbb C^{4\times4}\) on a carrier measure \(\mathrm d\mu_{\rm c}\) over the base, and let \(T_M\) be the operator \((T_M f)(x)=\int M(x,y)f(y)\,\mathrm d\mu_{\rm c}(y)\) on \(L^2(\mathrm d\mu_{\rm c})\otimes\mathbb C^4\), with \(T_W\) the operator of \(W\) alone and \(\kappa(T):=\sigma_{\max}(T)/\sigma_{\min}(T)\). (i)~Reduction. The factorization \(T_M=\mathsf P\,T_W\,\mathsf P^\dagger\), with \(\mathsf P\) fibrewise multiplication by \(\Phi\), applied to the identities \(\sigma_{\min}(AB)\ge\sigma_{\min}(A)\,\sigma_{\min}(B)\) and \(\sigma_{\max}(AB)\le\sigma_{\max}(A)\,\sigma_{\max}(B)\), gives: if \(c\le\sigma_{\min}(\Phi)\) and \(\sigma_{\max}(\Phi)\le C\) \(\mathrm d\mu_{\rm c}\)-a.e., then \(\sigma_{\min}(T_M)\ge c^2\sigma_{\min}(T_W)\) and \(\kappa(T_M)\le(C/c)^2\kappa(T_W)\). The gap recorded above therefore reduces exactly to the transport-kernel lower bound \(w:=\sigma_{\min}(T_W)\), and no bound on \(w\) follows from pointwise \(\Phi\)-data: with \(\Phi\equiv v\,\mathbb I\) and \(W\) fibrewise a rank-deficient projector, \(w=0=\sigma_{\min}(T_M)\) with \(\Phi\) invertible, and a fibre of smallest singular value \(c\) drives the bound at rate \(c^2\). (ii)~Unitary-transporter carrier case. If the carrier is finite and the transporter matrix is unitary -- matrix link variables and exactly unitary fuzzy Wilson lines, path-ordered exponentials of anti-Hermitian fibre connections, for which \(\sigma_{\min}=1\) exactly and uniformly in the carrier size -- then \(w=1\) and \(\sigma_{\min}(T_M)\ge c^2\); on a region where \(\mathcal V\le\varepsilon\) for \(\mathcal V\) of Equation~\eqref{eq:V}, the recorded bound \(\sigma_i(\Phi)^2\ge v^2-2\sqrt{\varepsilon/\lambda}\) upgrades this to \(\sigma_{\min}(T_M)\ge v^2-2\sqrt{\varepsilon/\lambda}>0\) for \(\varepsilon<\lambda v^4/4\). The estimate is an inequality in general, with equality on radial fibres \(\Phi=\rho\,U\), so it transfers non-degeneracy and does not determine the spectrum of \(T_M\). (iii)~Coincidence fibre and positivity. At coincidence, \(\sigma_{\min}(\Phi\Phi^\dagger)=\sigma_{\min}(\Phi)^2\) and \(\det(\Phi\Psi^\dagger)=\det\Phi\,\overline{\det\Psi}\) hold exactly, so the diagonal composite inherits pointwise non-degeneracy; every orthogonal compression of a positive-semidefinite fibre bounded below by \(c\) retains lower bound \(c\), while compression of a non-positive factor can drop below \(\sigma_{\min}\) -- positivity of the composite is what transports. (iv)~Diagonal state bound. For any one-point state \(\omega\) with potential moment \(\mu_{\mathcal V}:=\mathbb E_\omega[\mathcal V(x)]<\lambda v^4/4\), the Loewner estimate \(\lambda_{\min}(\mathbb E_\omega[\Phi\Phi^\dagger])\ge\mathbb E_\omega[\lambda_{\min}(\Phi\Phi^\dagger)]\), the fibre identity of (iii), and Markov's inequality give \(\lambda_{\min}(\mathbb E_\omega[\Phi\Phi^\dagger])\ge\sup_{\mu_{\mathcal V}<\varepsilon<\lambda v^4/4}\,(v^2-2\sqrt{\varepsilon/\lambda})(1-\mu_{\mathcal V}/\varepsilon)>0\): under the single moment hypothesis \(\mu_{\mathcal V}<\lambda v^4/4\) -- the one-point moment class made finite by the Gibbs paragraph above, whose smallness remains an input -- the coincidence composite \(\mathbb E_\omega[\Phi\Phi^\dagger]\) is state-nondegenerate; this is a statement about the composite and not about the mean-field frame \(e=\langle\Phi\rangle/v\) of Equation~\eqref{eq:vev-is-frame}, whose control is the mean-field gap recorded above. The hypothesis is sharp, the bound being non-positive at \(\mu_{\mathcal V}=\lambda v^4/4\) on the threshold fibre \((0,v,v,v)\), and no analogous bound holds for the off-diagonal expectation composite, whose fibre-wise Haar average vanishes identically. (v)~Carrier topology. The lower bound \(w\) is a property of the carrier, not of the kernel alone: on a periodic grid of \(n\) points the Fourier-diagonalised heat kernel of diffusion parameter \(s\) has smallest eigenvalue \(\exp(-s^2 n^2/4)\to0\), while its compression to a fixed band \(|k|\le K\) has grid-independent smallest eigenvalue \(\exp(-s^2K^2)\); a continuum smooth-kernel reading of the composite therefore cannot carry a uniform \(w>0\), and the bilocal non-degeneracy of Assumption~\ref{ass:nondeg} must be stated on a finite carrier as in (ii), the finite-\(N\) matrix-theory regime, or through the coincidence fibre (iii). For a path-ordered transporter \(W=P\exp\int\omega\) one has \(\sigma_{\min}(W)\ge\exp(-\int\sigma_{\max}(\omega))\), with equality on scalar-gauge paths and \(\sigma_{\min}=1\) for anti-Hermitian \(\omega\). Scope: (i)-(iii) and (v) are identities and conditional estimates under the stated a.e.\ bounds with explicit sharpness examples, and (iv) is a conditional estimate under the stated moment hypothesis; the carrier measure, the connection class of \(W\), and the state \(\omega\) are inputs and are not derived from Equation~\eqref{eq:Smicro}, no measure or state is constructed here, the off-diagonal expectation composite is not bounded, the vacuum-manifold fibre measure of Equation~\eqref{eq:vacuum-manifold} is untouched, and the computation of \(w\) for the specific carrier of Problem~D1 remains open; Assumptions~\ref{ass:frame-lock} and~\ref{ass:nondeg} remain assumed and this paragraph closes none of Problem~D, Problem~D1, Problem~F or Problem~N, all of which remain OPEN and PROGRAMME-CONDITIONAL.

